\documentclass[12pt,a4paper]{article}

\usepackage{fontawesome5}
\providecommand{\faPencilAlt}{\faIcon{pencil-alt}}

\usepackage{styles/couleurs}
\usepackage{styles/mise_en_page}
\usepackage{styles/boites}
\usepackage{styles/maths}
\usepackage{styles/figures}

\usepackage{enumitem}
\usepackage{array}
\usepackage{tabularx}
\usepackage{amssymb}
\usepackage{pdflscape}
\usepackage{tikz}
\usepackage{xparse}

\setlist[itemize]{leftmargin=1.4em,itemsep=0.15em,topsep=0.2em}
\setlist[enumerate]{leftmargin=*,itemsep=0.2em,topsep=0.2em}

\newcommand{\lignepointillee}{\noindent\dotfill\par\vspace{0.48em}}

\newcommand{\zonelignes}[1]{%
  \begin{tcolorbox}[
    enhanced,
    colback=white,
    colframe=NeutreC!40,
    boxrule=0.5pt,
    arc=1.5mm,
    left=3mm,
    right=3mm,
    top=2mm,
    bottom=1mm,
    before skip=0pt,
    after skip=0pt
  ]
    \foreach \i in {1,...,#1}{\lignepointillee}
  \end{tcolorbox}
}

\newcounter{exohc}

\NewDocumentCommand{\exohc}{m +m}{%
  \refstepcounter{exohc}
  \par\medskip
  \noindent
  \begin{minipage}[t]{0.47\linewidth}
    \vspace{0pt}
    {\sffamily\bfseries\color{TitreC}$\square$ Exercice \theexohc}
    \vspace{0.2em}
    #2
  \end{minipage}
  \hfill
  \begin{minipage}[t]{0.50\linewidth}
    \vspace{0pt}
    \zonelignes{#1}
  \end{minipage}
  \par\medskip
}

\begin{document}
\begin{landscape}

\begin{center}
  {\sffamily\LARGE\bfseries\color{TitreC}Exercices hors classe}\\[0.2em]
  {\sffamily\large\bfseries\color{ActionC}
  Cosinus d'un angle aigu -- Calculer des longueurs}
\end{center}

\section*{A. Vocabulaire et formule}

\exohc{4}{
Dans le triangle \(ABC\) rectangle en \(B\) :

\begin{enumerate}[label=\alph*)]
  \item nomme l'hypoténuse ;
  \item nomme le côté adjacent à \(\widehat A\) ;
  \item nomme le côté adjacent à \(\widehat C\).
\end{enumerate}
}

\exohc{4}{
Dans le triangle \(DEF\) rectangle en \(E\), complète :

\[
\cos(\widehat D)=\ldots\div\ldots
\]

\[
\cos(\widehat F)=\ldots\div\ldots
\]
}

\exohc{4}{
Donne une valeur approchée au centième :

\[
\cos 25^\circ
\qquad
\cos 48^\circ
\qquad
\cos 72^\circ.
\]

Puis vérifie que \(\cos 60^\circ=0{,}5\).
}

\section*{B. Calculer le côté adjacent}

\exohc{6}{
Dans le triangle \(GHI\) rectangle en \(H\),

\[
GI=8\text{ cm}
\qquad
\widehat G=37^\circ.
\]

Calcule \(GH\), arrondi au dixième de centimètre.
}

\exohc{6}{
Dans le triangle \(JKL\) rectangle en \(K\),

\[
JL=15\text{ m}
\qquad
\widehat L=64^\circ.
\]

Calcule \(KL\), arrondi au dixième de mètre.
}

\exohc{6}{
Dans le triangle \(MNP\) rectangle en \(N\),

\[
MP=5{,}6\text{ km}
\qquad
\widehat M=18^\circ.
\]

Calcule \(MN\), arrondi au centième de kilomètre.
}

\section*{C. Calculer l'hypoténuse}

\exohc{6}{
Dans le triangle \(RST\) rectangle en \(S\),

\[
RS=6\text{ cm}
\qquad
\widehat R=42^\circ.
\]

Calcule \(RT\), arrondi au dixième de centimètre.
}

\exohc{6}{
Dans le triangle \(UVW\) rectangle en \(V\),

\[
VW=3{,}5\text{ m}
\qquad
\widehat W=70^\circ.
\]

Calcule \(UW\), arrondi au dixième de mètre.
}

\exohc{6}{
Dans le triangle \(XYZ\) rectangle en \(Y\),

\[
XY=12{,}4\text{ mm}
\qquad
\widehat X=53^\circ.
\]

Calcule \(XZ\), arrondi au millimètre.
}

\section*{D. Choisir l'opération et contrôler}

\exohc{6}{
Sans calculer d'abord, indique s'il faut multiplier ou diviser.

\begin{enumerate}[label=\alph*)]
  \item On connaît l'hypoténuse et on cherche le côté adjacent.
  \item On connaît le côté adjacent et on cherche l'hypoténuse.
  \item Explique le choix à l'aide de la relation
  \[
  \text{adjacent}=\text{hypoténuse}\times\cos(\text{angle}).
  \]
\end{enumerate}
}

\exohc{6}{
Un élève obtient \(11{,}8\) cm pour le côté adjacent
d'un triangle rectangle dont l'hypoténuse mesure \(9\) cm.

Sans refaire le calcul, explique pourquoi le résultat est impossible.
}

\section*{E. Petits problèmes}

\exohc{7}{
Une échelle fait un angle de \(67^\circ\) avec le sol.
Son pied est à \(1{,}8\) m du mur.

Calcule la longueur de l'échelle nécessaire,
arrondie au dixième de mètre.
}

\exohc{7}{
Un câble de \(24\) m est tendu entre le sommet d'un pylône
et un point du sol. Il forme un angle de \(62^\circ\) avec le sol.

Calcule la distance horizontale entre le pylône et le point d'ancrage,
arrondie au dixième de mètre.
}

\exohc{8}{
Une observatrice se trouve à \(20\) m d'un arbre.
Ses yeux sont à \(1{,}60\) m du sol.
La ligne de visée vers le sommet forme un angle de \(40^\circ\)
avec l'horizontale.

\begin{enumerate}[label=\alph*)]
  \item Calcule la longueur de la ligne de visée avec le cosinus.
  \item Calcule ensuite la différence de hauteur avec Pythagore.
  \item Déduis la hauteur de l'arbre.
\end{enumerate}
}

\end{landscape}
\end{document}
